\documentclass[12pt]{article}
\usepackage[margin=1in]{geometry}
\usepackage{amsmath,amsthm,amssymb,scrextend}
\usepackage{pgfplots}
\usepackage{array}
\usepackage{tikz}
\setlength{\parindent}{0pt}
\begin{document}
\newcommand{\ds}{\displaystyle}

\textbf{(i)}

$$f(x)=x(x-a)(x-b)(x-c)$$

Area $0$ to $a$: $A$

$$F(0)=0,\  F(a)=-A,\  F(b)=A,\  F(c)=0$$

Turning points: $\ds x=0,\ a,\ b,\ c$

$$
\renewcommand{\arraystretch}{1.4}
\begin{array}{c!{\vrule width 1.2pt}c|c|c|c|c|c|c}
x & -\infty & 0 & a & h & b & c & \infty\\
\hline
\rule[-4pt]{0pt}{14pt}
F(x) & -\infty & 0 & -A & 0 & A & 0 & \infty
\end{array}
$$

\begin{center}
\begin{tikzpicture}
\begin{axis}[
    axis lines=middle,
    axis line style={->},
    xmin=-0.7,xmax=4.8,
    ymin=-4.5,ymax=4.5,
    xtick=\empty,
    ytick=\empty,
    width=11cm,
    height=7cm,
    samples=400,
    clip=false
]
\addplot[black,thick,domain=-0.65:4.65] {0.2*x^2*(x-4)^2*(x-2)};

\addplot[only marks,mark=*,mark size=1.5pt] coordinates {
    (0,0)
    ({2-2/sqrt(5)},{0.2*(2-2/sqrt(5))^2*((2-2/sqrt(5))-4)^2*((2-2/sqrt(5))-2)})
    (2,0)
    ({2+2/sqrt(5)},{0.2*(2+2/sqrt(5))^2*((2+2/sqrt(5))-4)^2*((2+2/sqrt(5))-2)})
    (4,0)
};

\node at (axis cs:0,0) [above left] {$0$};
\node at (axis cs:{2-2/sqrt(5)},0) [above] {$a$};
\node at (axis cs:2,0) [below] {$h$};
\node at (axis cs:{2+2/sqrt(5)},0) [below] {$b$};
\node at (axis cs:4,0) [below] {$c$};

\node at (axis cs:4.8,0) [right] {$x$};
\node at (axis cs:0,4.5) [above] {$y$};
\end{axis}
\end{tikzpicture}
\end{center}



$f(x)$ is monic quartic, so $F(x)$'s leading coefficient is $\frac15$

$$F(0)=0,\  F'(0)=0, F(c)=0,\ F'(c)=0$$

So $0,c$ are repeated roots of $F(x)$

$$F(a)<0,\  F(b)>0$$

So there is another root $h$ between $a$ and $b$

$$F(x)=\frac15x^2(x-c)^2(x-h),\quad F(c-x)=\frac15x^2(x-c)^2(c-h-x)$$


\begin{align*}
F(x)+F(c-x)
&=\frac15x^2(x-c)^2(x-h)
+\frac15x^2(x-c)^2(c-h-x)\\
&=\frac15x^2(x-c)^2(c-2h)
\end{align*}

\bigskip\hrule\bigskip

\textbf{(ii)}

For $0\leq x\leq c$
$$\quad F(x)\geq -A, \ F(x)=-A\iff x=a,\ F(b)=A$$

$$\therefore F(b)+F(x)\geq0$$

$$F(b)+F(x)>0\quad (x\neq a)$$

$$F(b)+F(c-b)=\frac15b^2(c-b)^2(c-2h)$$

$$c-b=a\quad\text{or}\quad F(b)+F(c-b)>0$$

$$c-b=a\quad\text{or}\quad c>2h$$



Similarly,
$$F(a)+F(x)\leq0,\quad F(a)+F(x)<0\quad (x\neq b)$$

$$F(a)+F(c-a)=\frac15a^2(c-a)^2(c-2h)$$

$$c-a=b\quad \text{or}\quad c<2h$$

$$c>2h\text{ and }c<2h\text{ cannot both hold}$$

$$\therefore c=a+b$$

$$F(b)+F(a)=0$$

$$\frac15b^2(c-b)^2(c-2h)=0\implies c=2h$$

\bigskip\hrule\bigskip

\textbf{(iii)}

$$h=\frac c2$$

$$F(x)=\frac15x^2(x-c)^2\left(x-\frac c2\right)$$

\begin{align*}
f(x)
&=F'(x)\\
&=\frac15x(x-c)(5x^2-5cx+c^2)\\
&=x(x-c)\left(x^2-cx+\frac{c^2}{5}\right)\\
&=\frac15(5x^4-10cx^3+6c^2x^2-c^3x)
\end{align*}

$$f''(x)=\frac15(60x^2-60cx+12c^2)=12\left(x^2-cx+\frac{c^2}{5}\right)$$

$$\because f(x)=x(x-c)\left(x^2-cx+\frac{c^2}{5}\right)$$

$$\therefore \text{When }f''(x)=0\implies f(x)=0$$

The points of inflection on $y=f(x)$ lie on the $x$-axis


\bigskip
\hrule
\begin{center}
    This solution was submitted by \textbf{Jon}\\[4pt]
    Published and refined by \textbf{Jon}\\[8pt]
    \textit{Last updated: July 18, 2026}
\end{center}

\end{document}